% Generated draft: requires mathematical and source review.
\documentclass[11pt]{article}
\usepackage[margin=25mm]{geometry}
\usepackage{fontspec}
\setmainfont{Latin Modern Roman}
\usepackage{amsmath,amssymb,mathtools,unicode-math}
\setmathfont{Latin Modern Math}
\usepackage{enumitem,needspace,xurl,hyperref}
\setlength{\parindent}{0pt}
\setlength{\parskip}{5pt}
\setlength{\emergencystretch}{3em}
\newcommand{\N}{\mathbb{N}}
\newcommand{\Z}{\mathbb{Z}}
\newcommand{\Q}{\mathbb{Q}}
\newcommand{\R}{\mathbb{R}}
\newcommand{\C}{\mathbb{C}}
\newcommand{\F}{\mathbb{F}}
\newcommand{\A}{\mathbb{A}}
\newcommand{\PP}{\mathbb{P}}
\newcommand{\E}{\mathbb{E}}
\newcommand{\eps}{\varepsilon}
\newcommand{\dd}{\,\mathrm{d}}
\begin{document}
\section*{UnsolvedMath ID: 1200037}
\textbf{Source problem code:} catalogue entry
\subsection{Definitions and mathematical statement}
In a standard measured formulation, a bounded-degree graphing is a symmetric Borel graph $E\subseteq X\times X$ without loops on a standard probability space $(X,\mu)$, generated by measure-preserving partial isomorphisms, with vertex degrees at most $D<\infty$. For a measurable edge set use the edge-counting measure
\[
 \nu(F)=\frac12\int_X\#\{y:(x,y)\in F\}\,\dd\mu(x).
\]
For measurable $A\subseteq X$, let $\partial_E A$ denote the edges with one endpoint in $A$ and one outside it. A measured-expansion convention requires some $h>0$ such that $\nu(\partial_E A)\geq h\mu(A)$ whenever $0<\mu(A)\leq1/2$. A strong-ergodicity convention requires that $\nu(\partial_E A_j)\to0$ imply $\mu(A_j)(1-\mu(A_j))\to0$ for every sequence of measurable sets $A_j$.

Fix an integer $C\geq1$. For a measurable coloring $c:X\to\{1,\ldots,C\}$ define
\[
 \operatorname{err}(c)=\nu\bigl(\{(x,y)\in E:c(x)=c(y)\}\bigr).
\]
Under these conventions, the question asks whether, for every graphing satisfying either the expansion hypothesis or the strong-ergodicity hypothesis,
\[
 \bigl[\forall\epsilon>0\ \exists\text{ measurable }c_\epsilon:X\to\{1,\ldots,C\},\
       \operatorname{err}(c_\epsilon)<\epsilon\bigr]
\]
\[
 \Longrightarrow
 \bigl[\exists\text{ measurable }c:X\to\{1,\ldots,C\},\
       \operatorname{err}(c)=0\bigr].
\]
The number $C$ is fixed throughout; it cannot grow as the error decreases.

The source does not give precise conventions for expansion, strong ergodicity, small error, or ``properly colored.'' The display specifies a natural measurable, almost-everywhere edge formulation consistent with the supplied assessment. Requiring a Borel coloring at every vertex, or a proper coloring on all edges including a null exceptional set, is a distinct target and is not silently identified with zero error. Bounded degree makes the small-edge-error and small-measure-of-incident-bad-vertices versions comparable, but their exact source convention remains to be confirmed.
\subsection{Short English statement}
For an expanding or strongly ergodic graphing, does the existence of measurable $C$-colorings with errors tending to zero force an exact measurable $C$-coloring? The supplied assessment provisionally labels this OPEN-TRIAGE, mentions partial understanding for factor-of-i.i.d. and approximate coloring, and explicitly says its literature search reached a cap. It supplies no fully resolving result or specific partial-result citation. This provisional assessment does not establish a definitive current open status.
\subsection{Sources}
\begin{itemize}
\item UnsolvedMath permanent ID 1200037, source code catalogue entry, ``Some Questions --- Question 37.'' Registry curation and metadata: Ulam AI and UnsolvedMath Contributors, CC BY 4.0; updated 2026-08-06.
\item Abert, \emph{Some questions} (2010), Question 37: \url{https://www.renyi.hu/~abert/questions.pdf}. Supplied manual transcription accessed 2026-07-29. The metadata says the 2010 author list was inspected but current status was not independently verified; status/release and rights review are NEEDS\_REVIEW.
\item The supplied provisional research assessment names no further authoritative reference for the claimed partial understanding or status. No missing citation is invented.
\end{itemize}
\end{document}
